\section{Results}\label{sec:results}
\begin{table}[htbp]
\centering
\small
\caption{Frozen model comparison. Accuracy, macro-F1, and top-5 are percentages; MRR is a unitless rank score. Values for LTD and DF/VarCNN come from S6; CUMUL-style from S7; Tik-Tok from S8. NetCLR not yet evaluated.}
\label{tab:main}
\setlength{\tabcolsep}{4pt}
\begin{tabular}{llrrrr}
\toprule
Cohort & Model & Acc. & Macro-F1 & Top-5 & MRR \\
\midrule
09B & LTD-Micro & 95.002 & 94.818 & 99.020 & 0.9685 \\
 & LTD-Macro & 84.068 & 83.946 & 98.523 & 0.9039 \\
 & \textbf{LTD-Hybrid} & \textbf{97.440} & \textbf{97.379} & \textbf{99.816} & \textbf{0.9853} \\
 & DF & 93.423 & 93.239 & 99.163 & 0.9593 \\
 & VarCNN & 97.290 & 97.216 & 99.775 & 0.9840 \\
 & CUMUL-style & 69.994 & 68.944 & 93.253 & 0.8023 \\
 & Tik-Tok & 94.110 & 93.858 & 99.135 & 0.9632 \\
\midrule
Future-B & LTD-Micro & 43.116 & 44.134 & 72.399 & 0.5615 \\
 & LTD-Macro & 22.305 & 18.328 & 57.995 & 0.3865 \\
 & \textbf{LTD-Hybrid} & \textbf{49.453} & \textbf{49.262} & \textbf{80.559} & \textbf{0.6306} \\
 & DF & 37.464 & 36.192 & 75.020 & 0.5358 \\
 & VarCNN & 33.447 & 31.140 & 67.244 & 0.4880 \\
 & CUMUL-style & 7.733 & 2.849 & 33.646 & 0.2265 \\
 & Tik-Tok & 20.105 & 18.557 & 52.214 & 0.3535 \\
\bottomrule
\end{tabular}
\end{table}

\paragraph{Temporally proximate evaluation (09B).} The hybrid attains 97.440\% accuracy, compared with 95.002\% for Micro and 84.068\% for Macro. The strongest SOTA comparator in this cohort is VarCNN at 97.290\%, a difference of only $0.150\pp$ from the hybrid. In the original S6 family, the hybrid--VarCNN accuracy contrast is not statistically distinguishable under the reported Holm adjustment; we therefore do not characterize the 09B result as confirmed superiority over VarCNN.

\paragraph{Later evaluation (Future-B).} All evaluated models undergo a substantial decrease in top-1 accuracy relative to 09B. The hybrid reaches 49.453\% accuracy and 49.262\% macro-F1 on Future-B, compared with Micro (43.116\%, 44.134\%), Macro (22.305\%, 18.328\%), DF (37.464\%, 36.192\%), and VarCNN (33.447\%, 31.140\%). The CUMUL-style SVM and Tik-Tok achieve 7.733\% and 20.105\% accuracy, respectively. These results are descriptive comparisons until the expanded paired-inference family is completed.

\paragraph{Temporal deterioration.} The hybrid's observed accuracy declines by 47.987 percentage points from 09B to Future-B; the corresponding differences for Micro, Macro, DF, and VarCNN are 51.886, 61.763, 55.959, and 63.843 points. The cohorts consist of different captures, so these differences characterize cross-cohort changes rather than paired changes for individual visits. This is consistent with a serious temporal generalization problem, but cannot by itself partition the effect among website evolution, software environments, and measurement conditions.

\paragraph{Incomplete comparisons.} No NetCLR accuracy, paired p-value for CUMUL/Tik-Tok, or Dataset C result is reported in this version. This omission is deliberate: the scientific narrative must not depend on anticipated findings.
